% !TEX root = ../main.tex
% Apercu de la charte et des schemas (a valider avant d'ecrire le plan complet)

% ======================================================================= titre
\begin{frame}[plain]
  \titlepage
\end{frame}

% ============================================================= ce que fait le moteur
\begin{frame}{Ce que fait le moteur}
  \framesubtitle{des particules, des ressorts, des forces}
  \begin{columns}[c,onlytextwidth]
    \begin{column}{0.58\textwidth}
      \begin{diagramoverlay}[width=\linewidth]{figures/scenes/scene_05.png}
        \astpoint{0.737}{0.94}{0.79}{0.80}{west}{ancre\\épinglée}
        \astpoint{0.722}{0.78}{0.79}{0.60}{west}{ressort\\étiré}
        \astpoint{0.5}{0.32}{0.5}{0.10}{north}{400 particules\astsep 1\,482 ressorts}
      \end{diagramoverlay}
    \end{column}
    \begin{column}{0.38\textwidth}
      \astbig{13}{scènes, du pendule au tissu}
      \vspace{1mm}
      \astbig[phelmared]{4}{forces : poids, frottement, ressort, vent}
      \vspace{1mm}
      \astbig[timagreen]{1/120\,s}{pas fixe, Euler ou Verlet en direct}
    \end{column}
  \end{columns}
  \astcue{
    \cue{montrer le tissu en direct, \cuekey{attraper une particule}}
    \cuenext{comment c'est organisé}
  }
\end{frame}

\section{Architecture}

% ================================================================ architecture
\begin{frame}{Quatre dossiers, une frontière}
  \framesubtitle{le moteur ne connaît que l'interface \code{Renderer}}
  \centering
  \begin{tikzpicture}[x=1cm,y=1cm]
    % physics : la bibliotheque
    \node[draw=floorgrey,line width=0.8pt,dash pattern=on 3pt off 2pt,rounded corners=3mm,
      fill=floorfill,minimum width=9.2cm,minimum height=3.3cm] (phys) at (4.6,1.65) {};
    \node[fill=white,draw=combislate,line width=0.6pt,font=\scriptsize\bfseries,
      inner xsep=2mm,inner ysep=1mm] at (phys.south) {PHYSICS/ \astsep{} LIBPHYSICS.A};
    \node[cls=datagreen/datafill,minimum width=1.9cm] (w) at (1.5,2.35) {\clsname{World}\clsbody{possède tout}};
    \node[cls=datagreen/datafill,minimum width=1.9cm] (p) at (3.7,2.35) {\clsname{Particle}\clsbody{position, masse}};
    \node[cls=datagreen/datafill,minimum width=1.9cm] (f) at (5.9,2.35) {\clsname{forces/}\clsbody{poids, ressort, vent}};
    \node[cls=datagreen/datafill,minimum width=1.9cm] (i) at (8.1,2.35) {\clsname{integrators/}\clsbody{Euler, Verlet}};
    \node[cls=memborder/memfill,abstract,minimum width=4cm] (r) at (4.6,0.95)
      {\clsabstract{Renderer}\clsbody{drawCircle, drawLine : virtuelles pures}};
    % scenes et app au-dessus
    \node[cls=dummyorange/dummyfill,minimum width=3.2cm] (sc) at (2.4,4.35)
      {\clsname{scenes/}\clsbody{13 scènes = 13 fonctions}};
    \node[cls=ctrlred/ctrlfill,minimum width=3.2cm] (app) at (6.8,4.35)
      {\clsname{app/main.cpp}\clsbody{boucle, clavier, souris}};
    % render : boite noire
    \node[blackbox,minimum width=3.0cm,minimum height=1.3cm] (sf) at (12.0,0.95)
      {{\bfseries\small render/}\\[0.3ex]SFMLRender\\{\tiny boîte noire}};
    \node[blackbox,minimum width=3.0cm] (lib) at (12.0,2.75) {SFML 2.6\\{\tiny FetchContent, statique}};
    % liens
    \draw[uses] (sc.south) -- (sc.south |- phys.north);
    \draw[uses] (app.south) -- (app.south |- phys.north);
    \draw[uses] (app.west) -- (sc.east);
    \draw[inherit] (sf.west) -- node[above,font=\tiny,text=textgrey]{implémente} (r.east);
    \draw[uses] (sf.north) -- (lib.south);
    \draw[uses] (app.east) -| (sf.north east);
    % la frontiere
    \draw[phelmared,line width=1.2pt,dash pattern=on 5pt off 3pt] (9.95,-0.25) -- (9.95,5.0);
    \node[font=\tiny\bfseries,text=phelmared,anchor=north,align=center] at (9.95,-0.3)
      {aucun en-tête SFML à gauche};
  \end{tikzpicture}
  \astcue{
    \cue{\cuekey{physics} compile en \code{.a}, sans SFML : vérifié avec \code{grep}}
    \cue{le rendu SFML : on n'en parle pas, c'est une boîte noire}
    \cuenext{ce qu'il y a dans physics}
  }
\end{frame}

% ================================================================ hierarchie
\begin{frame}{La hiérarchie des classes}
  \framesubtitle{une seule \code{Entity} partagée, grâce à l'héritage virtuel}
  \centering
  \begin{tikzpicture}[x=1cm,y=1cm]
    % racine
    \node[cls=combislate/white,minimum width=2.6cm] (E) at (5.6,5.75)
      {\clsname{Entity}\clsbody{id, name\astsep \code{static} current\_ids}};
    % interfaces
    \node[cls=memborder/memfill,abstract,minimum width=2.4cm] (D) at (3.0,3.95)
      {\clsabstract{Drawable}\clsbody{couleurs\\\code{draw(Renderer\&) = 0}}};
    \node[cls=memborder/memfill,abstract,minimum width=2.4cm] (F) at (8.2,3.95)
      {\clsabstract{Force}\clsbody{\code{compute(const Particle\&) = 0}}};
    % concret
    \node[cls=datagreen/datafill,minimum width=2.2cm] (P) at (0.95,2.05)
      {\clsname{Particle}\clsbody{pos, speed, mass\\forces : \code{Force\&}}};
    \node[cls=datagreen/datafill,minimum width=2.2cm] (S) at (5.6,2.05)
      {\clsname{Spring}\clsbody{$L_0$, $k$, $\rho$\\2 \code{Particle*}}};
    \node[cls=datagreen/datafill] (G) at (7.6,2.05) {\clsleaf{Gravity}\clsbody{$m\vec g$}};
    \node[cls=datagreen/datafill] (Dr) at (8.75,2.05) {\clsleaf{Drag}\clsbody{$-b\vec v$}};
    \node[cls=datagreen/datafill] (Wi) at (9.95,2.05) {\clsleaf{Wind}\clsbody{lambda $w(t)$}};
    % integrateurs
    \node[cls=memborder/memfill,abstract,minimum width=2.4cm] (I) at (12.9,3.95)
      {\clsabstract{Integrator}\clsbody{\code{step(particles, dt) = 0}}};
    \node[cls=datagreen/datafill] (Eu) at (12.15,2.05) {\clsleaf{Euler}\clsbody{semi-implicite}};
    \node[cls=datagreen/datafill] (Ve) at (13.65,2.05) {\clsleaf{Verlet}\clsbody{garde $x_{n-1}$}};
    % monde
    \node[cls=ctrlred/ctrlfill,minimum width=3.4cm] (W) at (5.6,0.2)
      {\clsname{World}\clsbody{copie interdite : \code{= delete}}};
    % heritages virtuels (le diamant)
    \draw[vinherit] (P.north) |- (E.west);
    \draw[vinherit,-] (D.north) -- (D.north |- E.west);
    \draw[vinherit] (F.north) |- (E.east);
    \draw[inherit] ([xshift=-5mm]S.north) |- (D.east);
    \draw[inherit] ([xshift=5mm]S.north) |- (F.west);
    \node[font=\tiny\bfseries,text=phelmared] at (5.6,3.95) {diamant};
    % heritages simples
    \draw[inherit] (P.east) -| ([xshift=-3mm]D.south);
    \draw[inherit] (G.north) -- (G.north |- F.south);
    \draw[inherit] (Dr.north) -- (Dr.north |- F.south);
    \draw[inherit] (Wi.north) |- (F.east);
    \draw[inherit] (Eu.north) -- (Eu.north |- I.south);
    \draw[inherit] (Ve.north) -- (Ve.north |- I.south);
    % possession
    \draw[owns] (W.west) -| node[pos=0.8,left,font=\tiny,text=textgrey,align=right]{\code{deque}\\\code{<Particle>}} (P.south);
    \draw[owns] (W.north) -- node[right,font=\tiny,text=textgrey]{\code{unique\_ptr<Spring>}} (S.south);
    \draw[owns] (W.east) -- node[above,pos=0.45,font=\tiny,text=textgrey]{\code{unique\_ptr<Integrator>}} (10.95,0.2) |- (I.west);
    % legende
    \begin{scope}[shift={(11.6,0.5)},font=\tiny,text=textgrey]
      \draw[vinherit] (0,0.45) -- (0.6,0.45); \node[anchor=west] at (0.7,0.45) {héritage virtuel};
      \draw[inherit] (0,0.1) -- (0.6,0.1); \node[anchor=west] at (0.7,0.1) {héritage};
      \draw[owns] (0,-0.25) -- (0.6,-0.25); \node[anchor=west] at (0.7,-0.25) {possède};
    \end{scope}
  \end{tikzpicture}
  \astcue{
    \cue{\cuekey{Spring} est à la fois une force et un dessin}
    \cue{sans \code{virtual} : deux \code{Entity}, deux identifiants}
    \cuenext{comment ces classes travaillent ensemble à chaque pas}
  }
\end{frame}

% ================================================================ boucle
\begin{frame}{Un pas de simulation}
  \framesubtitle{pas fixe : deux pas de physique par image}
  \centering
  \begin{tikzpicture}[x=1cm,y=1cm]
    % horloge
    \node[cls=combislate/white,minimum width=2.2cm,minimum height=1.25cm] (clk) at (0.9,2.6)
      {\clsname{Horloge}\clsbody{temps écoulé\\$\to$ accumulateur}};
    % la boucle de physique (etage)
    \node[draw=floorgrey,line width=0.8pt,dash pattern=on 3pt off 2pt,rounded corners=3mm,
      fill=floorfill,minimum width=8.6cm,minimum height=2.7cm] (loop) at (7.0,2.6) {};
    \node[fill=white,draw=combislate,line width=0.6pt,font=\scriptsize\bfseries,
      inner xsep=2mm,inner ysep=1mm] at (loop.south) {TANT QUE ACC $\geq$ DT};
    \node[cls=memborder/memfill,minimum width=2.3cm,minimum height=1.5cm] (s1) at (4.15,2.75)
      {\astpill{1}\\[0.4ex]\clsname{Forces}\clsbody{\code{f.compute(p)}\\virtuelle, par particule}};
    \node[cls=memborder/memfill,minimum width=2.3cm,minimum height=1.5cm] (s2) at (7.0,2.75)
      {\astpill{2}\\[0.4ex]\clsname{Intégration}\clsbody{\code{integrator->step()}\\Euler ou Verlet}};
    \node[cls=memborder/memfill,minimum width=2.3cm,minimum height=1.5cm] (s3) at (9.85,2.75)
      {\astpill{3}\\[0.4ex]\clsname{Contraintes}\clsbody{bords, rebond $e$\\points épinglés}};
    \draw[flow] (s1) -- (s2);
    \draw[flow] (s2) -- (s3);
    \draw[flow,rounded corners=3mm] (s3.north) -- ++(0,0.45) -| node[pos=0.25,above,font=\tiny,text=textgrey]{acc $-$= dt} (s1.north);
    \draw[flow] (clk) -- (loop.west);
    % dessin
    \node[cls=datagreen/datafill,minimum width=2.4cm,minimum height=1.25cm] (draw) at (13.1,2.6)
      {\clsname{Dessin}\clsbody{\code{d.draw(renderer)}\\pour chaque \code{Drawable}}};
    \draw[flow] (loop.east) -- (draw);
    \node[blackbox,minimum width=2.4cm] (sf) at (13.1,0.6) {SFML\\{\tiny boîte noire}};
    \draw[uses] (draw) -- (sf);
    % ce que World ne sait pas
    \node[font=\scriptsize,text=inpnavy,align=center] at (7.0,0.45)
      {\textbf{World} ne teste jamais un type : aucun \code{if}, aucun \code{switch}\\[0.3ex]
       {\color{textgrey}\tiny l'intégrateur se change en direct (touche I) : \code{enum class} + fabrique \code{makeIntegrator()}}};
  \end{tikzpicture}
  \astcue{
    \cue{pas fixe : \cuekey{Verlet suppose un dt constant}}
    \cue{1 \astsep{} 2 \astsep{} 3 = trois appels virtuels}
    \cuenext{les notions du cours, une par une}
  }
\end{frame}

\section{Les notions du cours}

% ================================================================ escalier
\begin{frame}{Du plus simple au plus avancé}
  \framesubtitle{chaque notion du cours, à sa place dans le code}
  \centering
  \begin{tikzpicture}[x=1cm,y=1cm]
    \foreach \i/\c/\f/\t/\b/\file in {
      0/datagreen/datafill/Vec2/{surcharge d'opérateurs\\\code{const\&}, méthodes \code{const}}/vec2.hpp,
      1/timagreen/timagreen!12/Constructeurs/{listes d'initialisation\\membre \code{const radius}\\compteur \code{static}}/entity.cpp,
      2/memborder/memfill/Polymorphisme/{classes abstraites\\\code{virtual}, \code{override}\\destructeur virtuel}/force.hpp,
      3/sinkpurple/sinkfill/Mémoire/{\code{unique\_ptr}, \code{deque}\\adresses stables\\règle des trois}/world.hpp,
      4/ctrlred/ctrlfill/Héritage multiple/{diamant virtuel\\\code{dynamic\_cast}\\lambdas, \code{std::function}}/springs.hpp}
    {
      \node[cls=\c/\f,anchor=south west,minimum width=2.75cm,minimum height=2.3cm,
        text width=2.5cm,align=left,inner sep=1.3mm] (st\i) at (\i*2.95,\i*0.85)
        {{\bfseries\footnotesize\t}\\[0.8ex]{\scriptsize\b}};
      \node[astnum] at (st\i.north west) {\the\numexpr\i+1\relax};
      \node[anchor=south east,inner sep=0.6mm] at (st\i.north east) {\filechip{\file}};
    }
    % la marche d'escalier, sous les boites
    \draw[-{Latex[length=2.4mm]},line width=1.4pt,phelmared,rounded corners=1mm]
      (0,-0.15) -- (2.85,-0.15) -- (2.85,0.7) -- (5.8,0.7) -- (5.8,1.55) -- (8.75,1.55)
      -- (8.75,2.4) -- (11.7,2.4) -- (11.7,3.25) -- (14.65,3.25) -- (14.65,4.6);
    \node[anchor=north east,font=\tiny\bfseries,text=phelmared,align=right] at (14.55,0.5)
      {plus abstrait,\\plus puissant};
  \end{tikzpicture}
  \astcue{
    \cue{on monte une marche par diapositive dans la suite}
    \cuenext{marche 1 : Vec2}
  }
\end{frame}

% ================================================================ dynamic_cast
\begin{frame}[fragile]{Attraper une particule à la souris}
  \framesubtitle{une lambda et un \code{dynamic\_cast}}
  \begin{columns}[c,onlytextwidth]
    \begin{column}{0.56\textwidth}
\begin{lstlisting}
Particle *World::pickParticle(const Vec2 &at,
                              float max_dist) {
  Particle *best = nullptr;
  float best_dist = max_dist;
  // la lambda capture best et best_dist
  auto closer = [&](Entity *e) {
    if (Particle *p = dynamic_cast<Particle *>(e)) {
      const float d = (p->getPos() - at).norm();
      if (d < best_dist) { best_dist = d; best = p; }
    }
  };
  std::vector<Entity *> all = entities();
  std::for_each(all.begin(), all.end(), closer);
  return best;
}
\end{lstlisting}
      \hfill\filechip{physics/world.cpp}
    \end{column}
    \begin{column}{0.40\textwidth}
      \centering
      \begin{tikzpicture}[x=1cm,y=1cm]
        \node[font=\scriptsize\bfseries,text=inpnavy] at (2.4,3.3) {\code{entities()} : des \code{Entity*}};
        \foreach \k/\t/\cc/\ff in {0/P/datagreen/datafill,1/S/memborder/memfill,2/P/datagreen/datafill,
                                   3/S/memborder/memfill,4/P/datagreen/datafill}
          \node[cls=\cc/\ff,minimum size=0.75cm,font=\small\bfseries] (e\k) at (\k*0.95+0.5,2.4) {\t};
        \foreach \k in {0,2,4} {
          \draw[flow,datagreen!70!black] (e\k.south) -- ++(0,-0.75) node[below,font=\tiny,text=datagreen!50!black]{\code{Particle*}};
        }
        \foreach \k in {1,3} {
          \draw[flow,floorgrey] (e\k.south) -- ++(0,-0.75) node[below,font=\tiny,text=floorgrey]{\code{nullptr}};
        }
        \node[astcallout,align=center,font=\scriptsize] at (2.4,0.1)
          {la plus proche du clic\\est attrapée};
      \end{tikzpicture}
    \end{column}
  \end{columns}
  \astcue{
    \cue{\code{dynamic\_cast} renvoie \code{nullptr} pour un ressort}
    \cue{la lambda capture par référence : \cuekey{[\&]}}
  }
\end{frame}

\section{Les scènes}

% ================================================================ galerie
\begin{frame}{Treize scènes, un seul moteur}
  \framesubtitle{une scène = une fonction qui remplit un \code{World}}
  \centering
  \setlength{\tabcolsep}{1.5mm}
  \begin{tabular}{ccc}
    \scenecard{figures/scenes/scene_04.png}{Pont de corde}{2.6cm} &
    \scenecard{figures/scenes/scene_05.png}{Tissu}{2.6cm} &
    \scenecard{figures/scenes/scene_07.png}{Rideau au vent}{2.6cm} \\[3mm]
    \scenecard{figures/scenes/scene_06.png}{Blob}{2.6cm} &
    \scenecard{figures/scenes/scene_12.png}{Toile d'araignée}{2.6cm} &
    \scenecard{figures/scenes/scene_10.png}{Pluie de balles}{2.6cm} \\
  \end{tabular}
  \astcue{
    \cue{ajouter une scène = écrire une fonction, \cuekey{sans toucher au moteur}}
  }
\end{frame}
